\subsection{The conic extension with the original parameter \texorpdfstring{$b$}{b} retained}
\label{sec:retained-b-conic}
All the constructions in this section are over $\mathbb C$. The
polynomial variables $b,s$ remain independent until a scalar value of
$b$ is explicitly fixed for the topological module construction.
The original coefficient is $a=2s$, and the original cubic on $c=0$ is
$-2r^2+br-4s$. Set
\[
 A=\mathbb C[b,s],\qquad
 B=A[r]/(r^2-br/2+2s),\qquad
 \beta=b/4,\quad \eta=r-\beta,
\]
and retain the inverse substitution $r=\eta+\beta$. Thus
\begin{equation}\label{eq:rb-original-base}
 s=br/4-r^2/2=b^2/32-\eta^2/2,\qquad
 \delta=b^2-32s=16\eta^2,\qquad
 L=4+b^2/16-2s=4+\eta^2.
\end{equation}
Here $\eta,\beta,\delta,L$ record explicit expressions in the original
objects; none changes the meanings of $b,r,s$ or $a$.

\subsubsection{The finite cover and its inverse maps}
Form the rank-two extension of $B$
\[
 E=B[p]/(p^2-L).
\]
Successive division by the two monic quadratic polynomials gives
the free $A$-basis $(1,r,p,rp)$: division first gives coefficients
of $1,p$ in $B$, and division in $B$ gives coefficients of $1,r$;
uniqueness at each step proves linear independence. There are inverse
$\mathbb C[b]$-algebra maps $E\longleftrightarrow\mathbb C[b,q,q^{-1}]$,
given on generators by
\begin{equation}\label{eq:rb-conic-maps}
 \begin{aligned}
 r&\longmapsto\beta+q-q^{-1},& p&\longmapsto q+q^{-1},\\
 s&\longmapsto b^2/32-(q-q^{-1})^2/2,& b&\longmapsto b,\\
 q&\longmapsto(p+r-\beta)/2,&q^{-1}&\longmapsto(p-r+\beta)/2.
 \end{aligned}
\end{equation}
To prove that these are well-defined, the two Laurent expressions
have $p^2-\eta^2=4$, and the proposed images of $q,q^{-1}$ have
product $(p^2-\eta^2)/4=1$. The image of $s$ then satisfies
\eqref{eq:rb-original-base}, hence both defining relations of $E$.
Conversely the sums and differences of the images of $q,q^{-1}$
recover $p,\eta$, and hence $r$ and $s$; the reverse compositions
recover $q,q^{-1}$ and $b$. This proves every map and its inverse.

For a temporary displayed coefficient put $h=2s-b^2/16$.
Then there are inverse $A$-algebra maps
\begin{equation}\label{eq:rb-quartic}
 E\simeq A[q]/\bigl(q^4+(2s-b^2/16-2)q^2+1\bigr),\qquad
 q^{-1}=-q^3+(2-2s+b^2/16)q.
\end{equation}
Indeed multiplying the proposed inverse of $q$ by $q$ gives one
by the quartic relation. In this quotient
$(q-q^{-1})^2=-h$, which follows by multiplying the quartic by
$q^{-2}$. Substituting in \eqref{eq:rb-conic-maps} therefore gives
the inverse map from $E$, and both compositions fix their generators.
In particular the quartic quotient is already a Laurent algebra.
Monic division gives its free basis $(1,q,q^2,q^3)$.

The complete conversion between these two ordered bases is
\begin{align}
 r&=\beta+q^3+(h-1)q,&
 p&=-q^3+(3-h)q,\notag\\
 rp&=h-2+\beta(3-h)q+2q^2-\beta q^3,&
 q&=(p+r-\beta)/2,\notag\\
 q^2&=(rp-\beta p+2-h)/2,&
 q^3&=r-\beta-(h-1)(p+r-\beta)/2.
 \label{eq:rb-basis-conversion}
\end{align}
The first two identities follow from the inverse of $q$ in
\eqref{eq:rb-quartic}. The identity
$\eta p=q^2-q^{-2}=2q^2+h-2$ proves the third and fifth.
The first identity proves the sixth after the fourth is inserted.
Thus no inverse basis map is omitted. The matrix whose columns are
$(1,r,p,rp)$ in the quartic basis is
\[
 T_b=\begin{pmatrix}
 1&\beta&0&h-2\\
 0&h-1&3-h&\beta(3-h)\\
 0&0&0&2\\
 0&1&-1&-\beta
 \end{pmatrix},\qquad \det T_b=4.
\]
Expanding along its first column and then the row containing $2$
gives $-2(-(h-1)-(3-h))=4$.

\subsubsection{Every original source coordinate on the cover}
Retain the original map, with its third source coordinate called $w$:
\begin{align*}
 F_1&=(1+xy)^3w+y^2(1+xy)(4+3xy),\\
 F_2&=y+3x(1+xy)^2w+3xy^2(4+3xy),\\
 F_3&=2x-3x^2y-x^3w.
\end{align*}
The simple-root inverse on $c=0$ is
\begin{equation}\label{eq:rb-source-H}
 \alpha=(b-4r)/2=-2\eta,\qquad
 H_b(\eta)=
 \left(-\frac1{2\eta},\frac b4+3\eta,
                 26\eta^2+\frac{3b\eta}{2}\right),\qquad\eta\ne0.
\end{equation}
These are the original formulas $x=1/\alpha$, $y=r-\alpha$,
$w=5\alpha^2-3r\alpha$ after the displayed invertible substitution.
Here $p$ is an additional coordinate, and never replaces $w$.
For a direct check of the full map set
$v=1+xy=-(\eta+\beta)/(2\eta)$ and note
$4+3xy=-(\eta+3\beta)/(2\eta)$. The polynomial identity
\[
 (\beta+3\eta)^2(\eta+3\beta)
 -(\eta+\beta)^2(13\eta+3\beta)=4\eta^2(\beta-\eta)
\]
follows by expanding both degree-three polynomials. Since
$w=2\eta(13\eta+3\beta)$, substitution in $F_1$ gives
$(\eta+\beta)(\beta-\eta)=\beta^2-\eta^2=2s$.
Substitution in $F_2$ gives
$\beta+3\eta+3(\beta-\eta)=4\beta=b$. Finally
\[
 F_3=-\frac1\eta-\frac{3(\beta+3\eta)}{4\eta^2}
                       +\frac{26\eta^2+6\beta\eta}{8\eta^3}=0.
\]
This proves $F(H_b(\eta))=(2s,b,0)$ with all coefficients retained.
Conversely the original inverse-chart coordinate is
$r=y+1/x$, so $\eta=y+1/x-b/4$; adjoining $p$ then recovers
$q=(p+y+1/x-b/4)/2$. The finite point
\[
 (x,y,w)=(0,b,2s-4b^2)
\]
lies on the same original fibre, because direct substitution gives
$F=(w+4y^2,y,0)=(2s,b,0)$. It remains separate from the $x\ne0$
chart, including at $s=b^2/32$.

The original inverse chart requires $\eta\ne0$, equivalently
$\delta\ne0$, or $q\ne\pm1$. The Laurent coordinate already
requires $q\ne0$. The extra quadratic extension has its additional
branch at $p=0$, equivalently $L=0$, or $q=\pm i$.
For completeness the relative differential module is exactly
\[
 \Omega_{E/A}=
 (E/(\eta))\,d\eta\ \oplus\ (E/(p))\,dp.
\]
Indeed $E=A[\eta,p]/(\eta^2-b^2/16+2s,p^2-L)$; the presentation
of relative differentials gives precisely the relations
$2\eta\,d\eta=0$ and $2p\,dp=0$, and $2$ is invertible.
Thus these are all the ramification loci, not just exhibited ones.
Their base values are
\[
 s=b^2/32\quad\hbox{and}\quad s=b^2/32+2,
\]
respectively. They are disjoint since their difference is $2$.
At the latter points $\eta=\pm2i$, so the original inverse chart
is still defined. At the former points $p=\pm2$; it is exactly the
original chart which excludes the point. These identities also
specify the domains of the derivative calculation below.

\subsubsection{Multiplication on the same arithmetic module}
Fix any scalar $b\in\mathbb C$, retaining its chosen value.
Let $Q$ be the original topological $\mathbb C[s]$-module with
continuous multiplication $S$ by its original coordinate $s$.
Give $Q$ the $A$-module structure $b\mapsto bI$, $s\mapsto S$ and
form $Q_b^{(4)}=E\otimes_AQ$. Since $E$ is free in the quartic basis,
the coefficient map gives an algebraic isomorphism $Q_b^{(4)}\simeq Q^4$;
we equip it with the fourfold direct-sum topology through this map.
With $\widehat S=\operatorname{diag}(S,S,S,S)$, multiplication by $q$ is
\begin{equation}\label{eq:rb-K}
 K_b=\begin{pmatrix}
 0&0&0&-I\\
 I&0&0&0\\
 0&I&0&(2+b^2/16)I-2S\\
 0&0&I&0
 \end{pmatrix}.
\end{equation}
The last column is the original quartic relation; the other three
columns shift the monic basis. Therefore
\[
 K_b^4+(2\widehat S-(b^2/16+2)I)K_b^2+I=0,\qquad
 K_b^{-1}=-K_b^3+((2+b^2/16)I-2\widehat S)K_b.
\]
Multiplication by this proposed inverse on either side gives $I$
by the quartic, because all its coefficients commute with $K_b$.
Both operators are continuous, being finite matrices of polynomials
in $S$. Define the original-root and added-coordinate operators by
\[
 \mathcal R_b=\frac b4I+K_b-K_b^{-1},\qquad
 \mathcal P_b=K_b+K_b^{-1}.
\]
The proved ring maps give
\begin{align}
 \mathcal R_b^2-\frac b2\mathcal R_b+2\widehat S&=0,&
 \mathcal P_b^2&=(4+b^2/16)I-2\widehat S,\notag\\
 \widehat S&=\frac b4\mathcal R_b-\frac12\mathcal R_b^2
 =\frac{b^2}{32}I-\frac12(K_b-K_b^{-1})^2.
 \label{eq:rb-original-S}
\end{align}
In the basis $(1,r,p,rp)$ multiplication by $r$ is
\[
 R_b\oplus R_b,\qquad
 R_b=\begin{pmatrix}0&-2S\\I&bI/2\end{pmatrix}.
\]
Indeed $r\cdot1=r$ and $r\cdot r=-2s+(b/2)r$; multiplying both
equations by $p$ gives the second block. The operator obtained by
substituting $S$ in $T_b$ is a continuous isomorphism: the inverse
formulas \eqref{eq:rb-basis-conversion} are finite polynomial
expressions in $S$ with scalar denominators $2$. Thus this is an
exact extension of the original two-sheet arithmetic module.

The continuous-invertibility spectrum of $\widehat S$ is precisely
that of $S$. If $S-\lambda I$ has a continuous two-sided inverse,
four copies give an inverse for $\widehat S-\lambda I$.
Conversely let $V$ be a continuous two-sided inverse of the latter.
For the first-coordinate inclusion $\iota:Q\to Q^4$ and projection
$\pi:Q^4\to Q$, the identities
$\widehat S\iota=\iota S$ and $\pi\widehat S=S\pi$ show that
$\pi V\iota$ is a continuous left and right inverse of
$S-\lambda I$. Also \eqref{eq:rb-original-S} gives the typed identity
\[
 S=\pi\left[\frac{b^2}{32}I
                 -\frac12(K_b-K_b^{-1})^2\right]\iota.
\]
Thus the recovered spectral coordinate is the original $s$.
An eigenvalue or spectral value of $K_b$ alone does not assert that
its recovered $s$ is off the original real axis.

Retain the GCT source-displayed Laurent coefficient polynomial
\[
 C(q)=q^{10}-q^4-q^{-4}+q^{-10}
      =(q-q^{-1})^2[7]_q[3]_q.
\]
The recurrence $U_0=1,U_1=p,U_j=pU_{j-1}-U_{j-2}$, after
$p=q+q^{-1}$, gives $U_j=q^j+q^{j-2}+\cdots+q^{-j}$ by cancellation
of the interior terms. Hence $[7]_q=p^6-5p^4+6p^2-1$ and
$[3]_q=p^2-1$. The factorization follows directly by multiplying
$(q-q^{-1})[7]_q=q^7-q^{-7}$ and
$(q-q^{-1})[3]_q=q^3-q^{-3}$; their product is exactly the
four-term Laurent expression above. Every factor of the coefficient
therefore descends by the explicit formula
\begin{align}
 C(q)&=(b^2/16-2s)(L^3-5L^2+6L-1)(L-1)\notag\\
 &=-(2s-b^2/16)\bigl(3-2s+b^2/16\bigr)\notag\\
 &\quad\cdot
 \bigl(7-14(2s-b^2/16)+7(2s-b^2/16)^2-(2s-b^2/16)^3\bigr).
 \label{eq:rb-C}
\end{align}
The second equality substitutes $L=4-(2s-b^2/16)$ and expands
only that polynomial; the first formula preserves its direct
conic substitution. In particular the coefficient has not become
the previous $b=0$ polynomial at the same $s$ unless $b=0$ is imposed.
For $\widehat L=(4+b^2/16)I-2\widehat S$, the actual transported
operator is exactly
\[
 C(K_b)=(b^2I/16-2\widehat S)
 (\widehat L^3-5\widehat L^2+6\widehat L-I)(\widehat L-I).
\]
Evaluation is legitimate by the explicit continuous inverse of $K_b$;
the identity follows from the proved Laurent algebra map. No factor
vanishing at a branch point has been divided out.

\subsubsection{The complete localized coefficient connection and its square}
Localize the base at both branch factors: $A^\circ=A[1/(\delta L)]$
and $E^\circ=E\otimes_AA^\circ$. The original material derivation on
the requested $c=0$ slice has $Db=0$, $Ds=1$, and consequently
$Da=2$. There is exactly one derivation of $E^\circ$ extending
these values, and it satisfies
\begin{equation}\label{eq:rb-derivation}
 D\eta=-1/\eta,\quad
 Dr=\frac4{b-4r}=\frac{4b-16r}{\delta},\quad
 Dp=-1/p,\quad
 Dq=-q/(\eta p).
\end{equation}
In fact differentiating the two monic quadratic relations forces
the first and third values because $\eta$ and $p$ are invertible
on this localization. Defining them this way annihilates both
relations, proving existence as well as uniqueness. Since
$Dr=D\eta$ and $b-4r=-4\eta$, the second formula follows, including
its expression in the free basis. In Laurent coordinates
$d\eta/dq=p/q$ and $dp/dq=\eta/q$; thus
\[
 D=-\frac{q}{\eta p}\frac{\partial}{\partial q}\quad(b\text{ fixed})
\]
has precisely these values and also gives $Ds=1$ in
\eqref{eq:rb-original-base}. Its Laurent domain is
$q\notin\{0,1,-1,i,-i\}$. This records separately the absence of
$q=0$, the original inverse-chart branch, and the added $p$ branch.

For a column $f=(f_0,f_1,f_2,f_3)^{\mathsf T}$ of coefficients in
$A^\circ$, interpret it as $f_0+rf_1+pf_2+rpf_3$.
The product rule and \eqref{eq:rb-derivation} give exactly
\begin{equation}\label{eq:rb-full-connection}
 Df=(\partial_sI_4+\Gamma_b)f,\qquad
 \Gamma_b=\begin{pmatrix}A_b&0\\0&A_b-I_2/L\end{pmatrix},\qquad
 A_b=\begin{pmatrix}0&4b/\delta\\0&-16/\delta\end{pmatrix}.
\end{equation}
Here the left side denotes the coefficient column of the derived
element, rather than componentwise differentiation alone.
The first two columns follow from $D1=0$ and
$Dr=(4b-16r)/\delta$. For the lower columns, $Dp=-p/L$ because
$p^2=L$, and
$D(rp)=p(4b-16r)/\delta-rp/L$. These verify all four columns.
Thus the $p$ contribution $-I_2/L$ is present in the full connection.

All zeroth-order terms of its square are
\begin{equation}\label{eq:rb-full-square}
 \begin{aligned}
 D^2&=\partial_s^2I_4+2\Gamma_b\partial_s+Z_b,\\
 Z_b&=\begin{pmatrix}
 (16/\delta)A_b&0\\
 0&(16/\delta-2/L)A_b-I_2/L^2
 \end{pmatrix}\\
 &=\begin{pmatrix}
 0&64b/\delta^2&0&0\\
 0&-256/\delta^2&0&0\\
 0&0&-1/L^2&64b/\delta^2-8b/(\delta L)\\
 0&0&0&-256/\delta^2+32/(\delta L)-1/L^2
 \end{pmatrix}.
 \end{aligned}
\end{equation}
For a full proof, $\partial_s\delta=-32$ and $\partial_sL=-2$ give
$A_b'=(32/\delta)A_b$ and $(-I_2/L)'=-2I_2/L^2$.
Direct matrix multiplication gives $A_b^2=-(16/\delta)A_b$.
It follows that the upper block of $\Gamma_b'+\Gamma_b^2$ is
$(16/\delta)A_b$. The lower block is
\[
 (32/\delta)A_b-2I_2/L^2
 -(16/\delta)A_b-2A_b/L+I_2/L^2
 =(16/\delta-2/L)A_b-I_2/L^2.
\]
Applying the product rule twice gives the displayed differential
operator and multiplying the scalars into $A_b$ gives every entry
of $Z_b$. In particular neither $-2A_b/L$ nor $-I_2/L^2$ disappears.

As a further compatibility check, the original root multiplication
obeys $R_b^2-(b/2)R_b+2sI_2=0$ and
\[
 [\partial_sI_2+A_b,R_b]
 =R_b'+A_bR_b-R_bA_b=4(bI_2-4R_b)/\delta.
\]
The last equality follows by substituting its four matrix entries:
the result is
$\delta^{-1}\left(\begin{smallmatrix}4b&32s\\-16&-4b\end{smallmatrix}\right)$.
It is exactly multiplication by $Dr$, establishing compatibility
with the original quadratic relation rather than only with its shift.

Restriction to $b=0$ is a quotient of this entire coefficient
calculation. Then $\eta=r$, $\delta=-32s$, $L=4-2s$, and the
connection becomes
\[
 \Gamma_0=\operatorname{diag}
 \left(0,\frac1{2s},-\frac1{4-2s},
                  \frac1{2s}-\frac1{4-2s}\right).
\]
Its full zeroth-order square term becomes
\[
 Z_0=\operatorname{diag}
 \left(0,-\frac1{4s^2},-\frac1{(4-2s)^2},
 -\frac1{4s^2}-\frac1{s(4-2s)}-\frac1{(4-2s)^2}\right).
\]
These follow by direct substitution in \eqref{eq:rb-full-square}.
The restricted domain is $s\ne0,2$, and the quartic, inverse source
map, original-root operator and coefficient in the preceding formulas
all specialize by setting $b=0$. Thus the earlier two-sheet connection
is exactly the first block of this four-sheet connection.

Finally the original coordinates in \eqref{eq:rb-source-H} have
the retained derivatives
\[
 Dx=-\frac1{2\eta^3},\qquad
 Dy=-\frac3\eta,\qquad
 Dw=-52-\frac{3b}{2\eta}.
\]
For example $Dw=(52\eta+3b/2)(-1/\eta)$; differentiation of the
other two displayed source functions proves their values.
Differentiating the already proved complete identity
$F(H_b(\eta))=(2s,b,0)$ now gives
\[
 DF(H_b(\eta))\,(Dx,Dy,Dw)^{\mathsf T}=(2,0,0)^{\mathsf T}.
\]
The separate finite point has derivative $(0,0,2)$.
For the original GCT coefficient, \eqref{eq:rb-C} gives the exact
identity $DC(q)=\partial_s C_b(s)$ with $b$ fixed, including all
its factors and the $b$-dependence. This derivation is on
$E^\circ$. The topological construction above proves multiplication
on $Q_b^{(4)}$ and its spectral-coordinate recovery; it does not
by itself define $\partial_s$ on $Q$, prove the invertibility there
of $\delta(S)$ or $L(S)$, or descend this localized connection to
an arithmetic quotient. No such derivative action is used in any
operator identity proved here.
